\chapter{Trabajo futuro}
\label{cap:futuro}

La evidencia actual se limita al simulador especificado de producto constante y sin comisiones. Los siguientes estudios abordan las brechas principales antes de formular afirmaciones sobre la operación de mercados.

\begin{enumerate}
  \item \textbf{Reproducir un evento de mercado documentado.} Seleccionar un evento con datos reproducibles y cuyo permiso permita redistribuirlos, como la desvinculación de USDC de marzo de 2023, y alinear los intercambios, reservas y eventos de liquidez DEX con registros de mejor compra/venta y profundidad CEX. Registrar el orden de bloque y transacción, las comisiones del pool, el gas, el retraso de ejecución, la orientación de los tokens, las restricciones de inventario y la procedencia de los datos. Fijar antes del análisis la ventana, la regla de sincronización y las métricas. La reproducción de un evento sería un estudio de caso, no una validación amplia; se requiere otro periodo o piscina para evaluar transferencia.
  \item \textbf{Rediseñar el experimento de respuesta de liquidez.} Asignar a los proveedores umbrales de retiro y fracciones retiradas heterogéneos, y comprobar que el disparador no se active en todas las corridas. Medir por separado la máxima desviación poschoque respecto de la nueva referencia y el sobrepaso direccional. Variar el tamaño de las órdenes relativo a las reservas y desacoplar la capacidad de arbitraje de las reservas de la piscina, para que el retiro no reduzca mecánicamente tanto la profundidad como el flujo de arbitraje. Usar un diseño factorial para estimar interacciones con el pánico e informar la incertidumbre Monte Carlo por semilla.
  \item \textbf{Incorporar costos de ejecución y competencia en cadena.} Implementar comisiones del pool, gas, retraso de liquidación, capital finito de los operadores y competencia entre buscadores de arbitraje. Relacionar estos mecanismos con los modelos de pérdida frente al rebalanceo y MEV antes de interpretar la capacidad como una limitación operativa.
  \item \textbf{Representar liquidez concentrada.} Extender la piscina de producto constante a posiciones delimitadas por rangos de precios, con liquidez activa que cambie al cruzar esos rangos. Conservar los resultados de producto constante como línea base deliberadamente acotada.
  \item \textbf{Evaluar modelos reducidos según su uso.} Mantener persistencia y AR(1) escalar como referencias; compararlos con un modelo sensible a la banda que use solo información disponible al momento de predecir. Evaluar en ventanas cronológicas y piscinas separadas. Si el objetivo es clasificar regímenes, definir etiquetas independientes de las variables del modelo y medir directamente el desempeño de clasificación.
  \item \textbf{Calibrar el tiempo del modelo.} Asociar los ticks con un reloj observable mediante marcas de tiempo o intervalos de bloque. Hasta hacerlo, la recuperación se mide en pasos de eventos y no puede interpretarse como segundos, bloques o latencia operativa.
  \item \textbf{Evaluar por separado ingeniería y pedagogía.} Comprobar si el tablero interactivo ayuda a ingenieros o estudiantes a distinguir el primer retorno a la banda del seguimiento persistente; esto evaluaría la interfaz, no validaría el modelo de mercado.
\end{enumerate}
